% Part: sets-functions-relations
% Chapter: sets
% Section: basics

\documentclass[../../../include/open-logic-section]{subfiles}

\begin{document}

\olfileid{sfr}{set}{bas}
\olsection{ઘટકો દ્વારા નિર્ધારિત સમાનતા}

\emph{ગણ} એટલે વસ્તુઓનો એવો સંગ્રહ, જેને {એક જ} વસ્તુ તરીકે ગણવામાં આવે.
ગણ બનાવતી વસ્તુઓને તેના \emph{ઘટકો} અથવા \emph{સભ્યો} કહે છે.
જો $x$ ગણ~$A$નો !!a{element} હોય, તો આપણે $x \in A$ લખીએ;
જો ન હોય, તો $x \notin A$ લખીએ. જે ગણમાં એક પણ !!{element} ન હોય
તેને \emph{ખાલી} ગણ કહે છે અને “$\emptyset$” વડે દર્શાવીએ છીએ.

\begin{explain}
આપણે ગણને કેવી રીતે \emph{નિર્દિષ્ટ} કરીએ છીએ, તેના !!{element}sને
કયા \emph{ક્રમમાં} મૂકીએ છીએ, કે તેના !!{element}sને \emph{કેટલી વાર}
ગણીએ છીએ, તેનાથી કોઈ ફરક પડતો નથી. તેના !!{element}s કયા છે એટલું જ
મહત્ત્વનું છે. આ વાતને આપણે નીચેના સિદ્ધાંતમાં વ્યક્ત કરીએ છીએ.
\end{explain}

\begin{defn}[ઘટકો દ્વારા નિર્ધારિત સમાનતા]
જો $A$ અને $B$ ગણો હોય, તો $A = B$ ત્યારે અને ત્યારે જ થાય જ્યારે
$A$નો દરેક !!{element} $B$નો પણ !!a{element} હોય અને ઊલટું પણ સાચું હોય.
\end{defn}

ઘટકો દ્વારા નિર્ધારિત સમાનતાનો સિદ્ધાંત કેટલીક સંજ્ઞાઓને યોગ્ય ઠરાવે છે.
સામાન્ય રીતે, આપણી પાસે $a_{1}$, \dots, $a_{n}$ વસ્તુઓ હોય, તો
$\{a_{1}, \dots, a_{n}\}$ એ \emph{ચોક્કસ એ જ} ગણ છે જેના !!{element}s
$a_1, \ldots, a_n$ છે. આપણે “\emph{ચોક્કસ એ જ}” શબ્દો પર ભાર મૂકીએ છીએ,
કારણ કે આ સિદ્ધાંત જણાવે છે કે આવો \emph{એક જ} ગણ હોઈ શકે.
હકીકતમાં, આ સિદ્ધાંત નીચેની સમાનતાઓને પણ યોગ્ય ઠરાવે છે:
  \[
    \{a, a, b\} = \{a, b\} = \{b,a\}.
  \]
આથી એ વાત સ્પષ્ટ થાય છે કે ગણોનો વિચાર કરતી વખતે તેમના !!{element}sનો
ક્રમ કે તેમને કેટલી વાર દર્શાવ્યા છે તે આપણા માટે મહત્ત્વનું નથી.

\begin{tagblock}{novice}
\begin{ex}
તમારી પાસે કેટલીક વસ્તુઓ હોય, ત્યારે તેમને એકત્ર કરીને ગણ બનાવી શકો છો.
દાખલા તરીકે, રિચર્ડનાં ભાઈબહેનોનો ગણ એક વ્યક્તિ ધરાવતો ગણ છે, અને તેને
$S=\{\textrm{રૂથ}\}$ તરીકે લખી શકાય.
$4$ કરતાં નાના ધન પૂર્ણાંકોનો ગણ $\{1, 2, 3\}$ છે, પરંતુ તેને
$\{3, 2, 1\}$ અથવા $\{1, 2, 1, 2, 3\}$ તરીકે પણ લખી શકાય.
ઘટકો દ્વારા નિર્ધારિત સમાનતાના સિદ્ધાંત અનુસાર આ બધા એક જ ગણ છે.
કારણ કે $\{1, 2, 3\}$નો દરેક !!{element} $\{3, 2, 1\}$નો
(અને $\{1, 2, 1, 2, 3\}$નો) પણ !!a{element} છે, અને ઊલટું પણ સાચું છે.
\end{ex}
\end{tagblock}

ઘણી વાર આપણે ગણના !!{element}sમાં રહેલા કોઈ સામાન્ય ગુણધર્મ દ્વારા
ગણને નિર્દિષ્ટ કરીશું. તે માટે આપણે આ સંક્ષિપ્ત સંજ્ઞા વાપરીશું:
$\Setabs{x}{\phi(x)}$. અહીં $\phi(x)$ એ એવો ગુણધર્મ દર્શાવે છે જે
ગણના !!{element}sમાં ગણાવા માટે $x$માં હોવો જરૂરી છે.

\begin{tagblock}{novice}
\begin{ex}
આપણા ઉદાહરણમાં $S$ને આ રીતે પણ નિર્દિષ્ટ કરી શક્યા હોત:
\[
S = \Setabs{x}{x \text{ રિચર્ડનાં ભાઈબહેનો પૈકી છે}}.
\]
\end{ex}
\end{tagblock}

\begin{tagblock}{math}
\begin{ex}
કોઈ સંખ્યાને \emph{પૂર્ણ સંખ્યા} ત્યારે અને ત્યારે જ કહે છે જ્યારે તે
તેના ઉચિત ભાજકોના સરવાળા જેટલી હોય (એટલે કે, તેને નિઃશેષ ભાગતા હોય
પણ તેના જેટલા ન હોય તેવા ભાજકો). દાખલા તરીકે, $6$ પૂર્ણ સંખ્યા છે,
કારણ કે તેના ઉચિત ભાજકો $1$, $2$ અને~$3$ છે, અને $6 = 1 + 2 + 3$.
હકીકતમાં, $10$ કરતાં નાના ધન પૂર્ણાંકોમાં $6$ એકમાત્ર પૂર્ણ સંખ્યા છે.
આથી, ઘટકો દ્વારા નિર્ધારિત સમાનતાનો સિદ્ધાંત વાપરીને કહી શકીએ:
\[
	\{6\} = \Setabs{x}{x\text{ પૂર્ણ સંખ્યા છે અને }0 \leq x \leq 10}
\]
જમણી બાજુની સંજ્ઞાને આપણે આ રીતે વાંચીએ: “એવા $x$નો ગણ કે $x$
પૂર્ણ સંખ્યા છે અને $0 \leq x \leq 10$”. અહીંની સમાનતા ખાતરી આપે છે
કે ગણોનો વિચાર કરતી વખતે તેમને કેવી રીતે નિર્દિષ્ટ કર્યા છે તે
મહત્ત્વનું નથી. વધુ સામાન્ય રીતે, ઘટકો દ્વારા નિર્ધારિત સમાનતાનો
સિદ્ધાંત ખાતરી આપે છે કે $\phi(x)$ હોય તેવા $x$નો હંમેશાં એક જ ગણ હોય.
આથી $\Setabs{x}{\phi(x)}$ને $\phi(x)$ હોય તેવા $x$નો
\emph{ચોક્કસ એ જ} ગણ કહેવું યોગ્ય છે.
\end{ex}
\end{tagblock}

ઘટકો દ્વારા નિર્ધારિત સમાનતાનો સિદ્ધાંત ગણો એક જ છે તે બતાવવાની રીત આપે છે:
$A = B$ બતાવવા માટે, જ્યારે પણ $x \in A$ હોય ત્યારે $x \in B$ પણ હોય,
અને જ્યારે પણ $y \in B$ હોય ત્યારે $y \in A$ પણ હોય, તે બતાવો.

\begin{prob}
વધુમાં વધુ એક જ ખાલી ગણ હોઈ શકે તે સાબિત કરો. એટલે કે, જો $A$ અને $B$
એવા ગણો હોય જેમાં કોઈ !!{element}s ન હોય, તો $A = B$ થાય તે બતાવો.
\end{prob}

\end{document}
